% ============================== 行间距与页边距 ==============================
\usepackage[top=2.5cm,bottom=2.5cm,left=2.5cm,right=2.5cm]{geometry}
\usepackage{setspace}
\setstretch{1.25} % 设置1.25倍行距

% ============================== 头文件与页眉页脚 ==============================
\usepackage{fancyhdr}
\usepackage{lastpage} % 用于获取总页数
\setlength{\headheight}{14.5pt} % 防止 fancyhdr 警告
\pagestyle{fancy}
\fancyhf{} % 清空页眉页脚
% 页眉：靠左显示实验编号与名称
\fancyhead[L]{\songti\zihao{5} 实 验 6 一 四 探 针 测 试 技 术}
% 页脚：居中显示 "第 X 页 共 Y 页"
\fancyfoot[C]{\songti\zihao{5} 第 \thepage\ 页 共 \pageref{LastPage} 页}
% 页眉分割线粗细
\renewcommand{\headrulewidth}{0.4pt}

% ============================== 章节标题格式重构 ==============================
\usepackage{titlesec}
\usepackage{enumitem}

% 一级标题：如 "一、实验目的"
\titleformat{\section}{\large\heiti\bfseries}{\chinese{section}、}{0em}{}
\titlespacing{\section}{0pt}{1.5ex plus .1ex minus .2ex}{1ex plus .1ex}

% 二级标题：如 "（一）分子能级与光谱"
\titleformat{\subsection}{\normalsize\heiti}{\thesubsection}{0em}{}
\renewcommand{\thesubsection}{（\chinese{subsection}）} % 重新定义计数器外观
\titlespacing{\subsection}{2em}{1ex plus .1ex minus .2ex}{0.5ex plus .1ex}

% ============================== 插图与表格标题格式 ==============================
\usepackage{caption}
\usepackage{graphicx}
\usepackage{amsmath}
\usepackage{amssymb} % 提供 \square 等符号
\usepackage{booktabs}
\usepackage{array}
\usepackage[hidelinks]{hyperref}
\captionsetup[figure]{labelsep=space, font=small}
\captionsetup[table]{labelsep=space, font=small}

% ============================== 其他常用宏包 ==============================
\usepackage{indentfirst} % 首行缩进
\usepackage{float}       % 图表浮动体控制 [H]
